%% =====================================================================
%% Discussion / Limitations and future work -- owner: Harshaan.
%% No HySortK here (Nov 13 decision). No O3/O5 labels: research-plan.md
%% and later docs disagree on what O3 means; O5 (GPU) is not discussed.
%% Nothing below is implemented.
%% =====================================================================

\subsection{Limitations and future work}\label{sec:disc:future}

Alignment is now the dominant stage, and its cost is unevenly distributed
across ranks. The time of an alignment task follows the amount of sequence
it aligns, which is not known in advance and is poorly predicted by seed
counts or contig lengths. This suggests replacing the static assignment of
contig pairs with dynamic scheduling, in which ranks take the next task from
a shared queue as they become free.
% CHECK: "dominant" + "poorly predicted" come from decision-gate-2
%        (primates); keep qualitative here, numbers in the evaluation.

Scheduling alone has a limit: each contig pair is aligned by one rank, and
self-alignments of long contigs are the largest tasks. Scaling further
requires splitting such a task across ranks while producing exactly
FastGA's alignments. This is an open problem.

Two implementation limits also bound scale. The pair-to-rank map uses 8-bit
entries, which caps the number of ranks at 127, and extraction partitions
work at whole-contig boundaries.

Finally, exact validation at human scale requires an independent
single-node FastGA reference, whose cost grows rapidly with the number of
haplotypes. Agreement between runs at different rank counts is weaker
evidence than agreement with FastGA, because it shares any error of our
own pipeline.
% TODO(Harshaan): Add human results and interpretation after validation
%   (branches A/B/C in disc-human.tex; checklist in author-notes).
